
We introduce \textbf{Doob Policy Optimization (DoobPO)}, a general framework for RL fine-tuning of diffusion policies (\cref{sec:diff_dr_matching}) and flow policies (\cref{sec:flow_dr_matching}) without evaluating exact action likelihoods.
DoobPO supports a broad family of trust-region policy gradient methods, while \cref{sec:unified_framework} shows how both existing and proposed algorithms arise under the same framework.

\subsection{Diffusion Policy Optimization with DoobPO}
\label{sec:diff_dr_matching}

In this section, we introduce DoobPO for diffusion policies.
We start from a base diffusion model with score network $s_\phi(x,t)$, which induces a diffusion policy $\pi_\phi$.
Our goal is to optimize the policy objective \eqref{eq:rl_objective_prelim} over parameters $\theta$, starting from the pretrained initialization $\theta=\phi$.

\textbf{Ratio Learning.}
We first learn a probability ratio that captures the policy improvement signal induced by trust-region policy-gradient updates, including regularized Policy Mirror Descent \cite{zhan2023policy}, Soft Actor-Critic \cite{haarnoja2018soft}, PPO \cite{schulman2017proximal}, and GRPO \cite{shao2024deepseekmath}.
Specifically, we introduce a lightweight ratio network $r_\beta(o,a)$ to approximate the intractable likelihood ratio
$r_\theta(o,a)=\pi_\theta(a|o)/\pi_{\mathrm{old}}(a|o)$.
We train $r_\beta$ on trajectories collected under the behavior policy $\pi_{\mathrm{old}}$ by optimizing a clipped reinforcement learning objective $J_{\mathrm{RL}}$.
For example, under PPO, we use
$J_{\mathrm{RL}} = J_{\mathrm{PPO}}$
following \eqref{eq:ppo}, with $r_\theta$ replaced by $r_\beta$.
Intuitively, $r_\beta(o,a)>1$ upweights an action relative to $\pi_{\mathrm{old}}$, while $r_\beta(o,a)<1$ downweights it.
The clipping term enforces a conservative trust-region-style update.
In practice, we parameterize the ratio model as
$r_\beta(o,a)=\exp(\Tilde{r}_\beta(o,a))$ to ensure $r_\beta(o,a)>0$.


\textbf{Doob's Drift Matching.}
We now distill the policy improvement information encoded by the learned ratio $r_\beta$ into the score model.
We consider a diffusion policy, where the denoising SDE \eqref{eq:reverse_sde} is conditioned on the environment state $o$, and the score estimator takes the form $s_\phi(x,t,o)$.
This score model induces the behavior policy $\pi_{\mathrm{old}}(\cdot|o)$ as the terminal distribution of the denoising process.

Given the learned ratio $r_\beta(o,a)$, we would like to update the policy toward the ratio-tilted distribution
$\pi(\cdot|o)\propto r_\beta(o,\cdot)\pi_{\mathrm{old}}(\cdot|o)$.
We realize this tilt by modifying the denoising dynamics with an additional drift term.
We characterize this modification via Doob's $h$-transform.

\begin{lemma}[Doob's $h$-Transform]
\label{lem:algo_correctness}
Fix a state $o \in\mathcal{S}$.
Let $(X_t)_{t\in[0,T]}$ follow the reverse-time denoising SDE \eqref{eq:reverse_sde} with $\nabla_x \log p_t(x|o)$ replaced by its estimator $s_\phi(x,t,o)$, inducing conditional marginals
$\{p_t^\phi(\cdot|o)\}_{t\in[0,T]}$.
Let $h_0(\cdot,o):\mathcal{A}\to\R_+$ be a positive integrable function.
Let $h(\cdot,o,\cdot):\mathcal{A}\times[0,T]\to \R_+$ be the associated Doob $h$-function, denoted by $h_t(x,o)$, defined as
\begin{align}
\label{eq:solution_exp}
    h_t(x,o)
    =
    \E \big[h_0(X_0,o)| X_t=x,o\big],
\end{align}
where the conditional expectation is taken under the original (untilted) denoising process induced by $s_\phi$.
The Doob $h$-transformed denoising process starts from
\[
    X_T^h \sim p_T^h(\cdot|o),
    \qquad
    p_T^h(x|o)
    :=
    \frac{
        p_T^\phi(x|o) h_T(x,o)
    }{
        \E_{X_0\sim p_0^{\phi}(\cdot|o)}[h_0(X_0,o)]
    },
\]
and evolves backward from $T$ to $0$ according to
\begin{align}
\label{eq:new_process}
    \dd X_t^{h}
    =
    \Big[
        f(X_t^{h},t)
        -\sigma_t^2 s_\phi(X_t^{h},t,o)
        -\sigma_t^2 \nabla_x \log h_t(X_t^{h},o)
    \Big]\dd t
    +\sigma_t\dd W_t.
\end{align}
Then its conditional marginals satisfy
\begin{align}
    p_t^{h}(x|o)
    =
    \frac{
        p_t^\phi(x|o)h_t(x,o)
    }{
        \E_{X_0\sim p_0^{\phi}(\cdot|o)}[h_0(X_0,o)]
    },
    \quad t\in[0,T].
    \label{eq:p_t^h}
\end{align}
Particularly, taking $h_0(x,o)=r_\beta(o,x)$ yields
$p_0^h(x|o)\propto p_0^{\phi}(x|o)r_\beta(o,x)$.
With $p_0^{\phi}(\cdot|o)=\pi_{\mathrm{old}}(\cdot|o)$, the time-$0$ policy induced by the tilted denoising process satisfies
$\pi(a|o)\propto r_\beta(o,a)\pi_{\mathrm{old}}(a|o)$.
\end{lemma}

\begin{proof}
Please see \cref{sec:proof_corr} for a detailed proof.
\end{proof}


\begin{remark}[Initialization]
In \cref{lem:algo_correctness}, the Doob-transformed denoising process \eqref{eq:new_process} is initialized from
$p_T^h(\cdot|o)$, which does not necessarily coincide with the original terminal law
$p_T^\phi(\cdot|o)$.
Under the idealized fully-noised setting, where $X_T$ is independent of $X_0$ conditional on $o$, we have
\[
    h_T(x,o)
    =
    \mathbb E[h_0(X_0,o)|X_T=x,o]
    =
    \mathbb E_{X_0\sim p_0^\phi(\cdot|o)}[h_0(X_0,o)].
\]
Therefore, $p_T^h(\cdot|o)=p_T^\phi(\cdot|o)$ under this assumption.
In practice, we initialize \eqref{eq:new_process} from the same standard Gaussian prior as the reference process.
\end{remark}


We then learn a parameterized conditional score model $s_\theta(x,t,o)$ to match the ratio-induced drift $\nabla_x \log h_t(x,o)$ via the Doob's Drift Matching objective
\begin{align}
\label{eq:def_DM}
    \mathcal{L}_{\mathrm{DM}}(\theta)
    =
    \E_{t, o,X_t \sim p_t^{h}(\cdot|o)}
    \Big[
        \big\|
            \nabla_x \log h_t(X_t,o)
            - s_\theta(X_t,t,o)
        \big\|_2^2
    \Big],
\end{align}
where $\theta$ denotes the learnable parameters and $t$ samples from uniform distribution $U[0,T]$.

Although the Doob's Drift Matching (DM) objective is clean and direct, it requires access to both $\nabla_x \log h_t(x,o)$
and samples from the tilted distribution $p_t^{h}(\cdot|o)$, which are generally intractable.
To address this issue, we introduce a conditional objective that attains the same optimum as $\mathcal{L}_{\mathrm{DM}}$.
Specifically, we define the Conditional Drift Matching (CDM) objective
\begin{align}
\label{eq:def_cDM}
    \mathcal{L}_{\mathrm{CDM}}(\theta)
    =
    \E_{t,o, X_t \sim p_t^\phi, X_0 \sim p_{0| t}^\phi}
    \Big[
        r_\beta(o,X_0)
        \big\|
            \nabla_x \log p_{0|t}^\phi (X_0| X_t,o)
            - s_\theta(X_t,t,o)
        \big\|_2^2
    \Big],
\end{align}
where $p_{0| t}^\phi(\cdot| X_t,o)$ denotes the diffusion conditional distribution.
More importantly, for a fixed state $o$, we show that the objectives in \eqref{eq:def_DM} and \eqref{eq:def_cDM} share the same gradient direction.

\begin{theorem}[Gradient Equivalence for Diffusion]
\label{thm:grad_equi_diff}
Fix $o \in\mathcal S$, define normalizing constant
$Z_{\beta}(o) := \E_{X_0 \sim p_0^\phi(\cdot|o)}\big[r_\beta(o,X_0)\big]$ and
\begin{align*}
    \ell_{\mathrm{DM}}(\theta,o)
    &:=
    \E_{t,X_t\sim p_t^h(\cdot|o)}
    \left[
        \|
        \nabla_x\log h_t(X_t,o)-s_\theta(X_t,t,o)
        \|_2^2
    \right],
    \\
    \ell_{\mathrm{CDM}}(\theta,o)
    &:=
    \E_{t,X_t\sim p_t^\phi(\cdot|o),X_0\sim p_{0|t}^\phi(\cdot|X_t,o)}
    \left[
        r_\beta(o,X_0)
        \|
        \nabla_x\log p_{0|t}^\phi(X_0|X_t,o)
        -
        s_\theta(X_t,t,o)
        \|_2^2
    \right].
\end{align*}
Then their gradients satisfy
$\nabla_\theta \ell_{\mathrm{CDM}}(\theta,o)
=
Z_{\beta}(o) \cdot \nabla_\theta \ell_{\mathrm{DM}}(\theta,o)$.
\end{theorem}

\begin{proof}
Please see \cref{sec:proof_grad_equi} for a detailed proof.
\end{proof}

\cref{thm:grad_equi_diff} is stated for each fixed observation $o$.
After averaging over observations, we obtain
\begin{align*}
    \nabla_\theta \mathcal L_{\mathrm{CDM}}(\theta)
    =
    \nabla_\theta \E_o \left[Z_{\beta}(o)\ell_{\mathrm{DM}}(\theta,o)\right].
\end{align*}
Therefore, $\mathcal L_{\mathrm{CDM}}$ has the same gradient field as the reweighted DM objective
$\E_o[Z_{\beta}(o)\ell_{\mathrm{DM}}(\theta,o)]$.
If the model class $\{s_\theta:\theta\in\Theta\}$ is realizable and the minimum of $\mathcal L_{\mathrm{CDM}}$ is attained, this reweighting does not change the global minimizer.
Indeed, there exists $\theta^\star\in\Theta$ such that
$s_{\theta^\star}(x,t,o)=\nabla_x\log h_t(x,o)$
for every $x\in\mathcal A$, $t\in[0,T]$, and $o\in\mathcal S$.
Hence $\ell_{\mathrm{DM}}(\theta^\star,o)=0$ for every $o\in\mathcal S$,
so $\theta^\star$ also minimizes both the original DM objective and the reweighted DM objective, which implies
$\theta^\star$ also minimizes $\mathcal L_{\mathrm{CDM}}$.


For variance-preserving DDPMs, the forward noising
process \eqref{eq:forward_sde} has the closed-form
representation
$X_t = \sqrt{\bar\alpha_t}X_0
    + \sqrt{1-\bar\alpha_t}Z,
    \quad Z\sim\mathcal N(0,I)$,
where
$\bar\alpha_t
= \exp(-\int_0^t \sigma_s^2 \dd s)$
is the cumulative signal retention coefficient.
Moreover, the inference process follows the backward SDE \eqref{eq:reverse_sde} by replacing the drift term by $-\sigma_t^2 x/2 - \sigma_t^2 s_\phi(x,t)$.
Therefore, following the Bayes formula $p_{0|t} = p_{t|0} p_0/p_t$, the CDM objective takes the form
\begin{align}
    \mathcal{L}_{\text{CDM}}(\theta)
    = & ~ \E_{t,o, X_t, X_0}
    [r_\beta(o,X_0) \|\nabla \log p_{t|0}^\phi(X_t|X_0,o) - (s_\phi + s_\theta)(X_t,t,o) \|_2^2] \nonumber \\
    \approx & ~\E_{t,o, X_t, X_0}
    [ r_\beta(o,X_0) \| -\frac{X_t - \sqrt{\bar{\alpha}_t}X_0}{1-\bar{\alpha}_t}- (s_\phi + s_\theta)(X_t,t,o) \|_2^2],
    \label{eq:cdm_loss_ddpm}
\end{align}
where $X_t \sim p_t$, $t \sim U[0,T]$, and $X_0$ is sampled from the conditional distribution
$p_{0|t}(\cdot| X_t)$.
Let the score model $s_\phi$ be the pre-trained base model that fixes the noise schedule and the corresponding regression target.
The approximation in the second line replaces the conditional score
$\nabla \log p_{t|0}^\phi(X_t|X_0,o)$ by the Gaussian regression target induced by the forward noising kernel.
The approximation is reliable when $s_\phi$ is well-trained under the chosen diffusion construction and the trust-region update keeps $\pi_{\mathrm{old}}$ close to the base behavior.

\textbf{Drift Parametrization.}
The conditional drift matching loss \eqref{eq:cdm_loss_ddpm} requires evaluating
the base score network $s_\phi$.
In the advantage-tilted denoising process \eqref{eq:new_process}, both the base
score $s_\phi$ and the advantage-induced drift $\nabla_x \log h_t$ appear with
the same coefficient $\sigma_t^2$.
We exploit this structure and adopt a simplified parametrization.
Specifically, we parametrize the drift model $g_\theta$ to match the sum
$\nabla_x \log h_t(x) + s_\phi(x,t)$.
Then the conditional drift matching loss
\eqref{eq:cdm_loss_ddpm} reduces to
\begin{align}
    \mathcal{L}_{CDM}(\theta)
    = \E_{t,o, X_t, X_0}
    \Big[r_\beta(o,X_0)
    \| -\frac{X_t - \sqrt{\bar{\alpha}_t}X_0}{1-\bar{\alpha}_t} - g_\theta(X_t,t,o) \|_2^2 \Big],
    \label{eq:cdm_loss_prac}
\end{align}
which no longer requires evaluating the base score network $s_\phi$ during
training.
This parametrization also simplifies inference.
Once we obtain the trained score network $g_\theta$, sampling from the target distribution $p_0^h$ only requires simulating the
following SDE:
\begin{align}
\label{eq:new_process_prac}
    \dd X_t^h
    = [f(X_t^h,t)
    - \sigma_t^2 g_\theta(X_t^h,t,o)
    ] \dd t + \sigma_t \dd W_t.
\end{align}
We summarize the training procedure in \cref{alg:rwdm_ppo}.

\subsection{Flow Policy Optimization with DoobPO}
\label{sec:flow_dr_matching}

In this section, we instantiate DoobPO for flow policies.
We start from a pre-trained flow matching model $u_\phi(x,t,o)$, which induces a flow policy $\pi_\phi(\cdot|o)$ with ODE \eqref{eq:fm_ode}.
Our goal is to solve \eqref{eq:rl_objective_prelim}.
We reuse the same ratio learning stage as in \cref{sec:diff_dr_matching}.

\textbf{Doob's Drift Matching.}
We then distill the policy improvement signal encoded by $r_\beta$ into the velocity model.
For flow matching, the policy is given by the terminal distribution at time $T$.
We consider the joint distribution $p_{t,T}^\phi$
induced by the independent-coupling affine interpolation
$X_t=\alpha_t X_0+\beta_t X_T$, where
$X_0\sim\mathcal N(0,I)$,
$X_T\sim\pi_{\mathrm{old}}(\cdot|o)$,
and $X_0\perp X_T$.
We assume the pretrained flow matching model
learns the marginal velocity field of
this interpolation.
We define the tilted flow path by terminally reweighting the base joint law:
$p_{t,T}^\beta(x,a|o)
:=
    {
        r_\beta(o,a)
        p_{t,T}^\phi(x,a|o)
    }/{
        Z_{\beta}(o)
    },
    ~~ \text{where} ~~
    Z_{\beta}(o):=
    \E_{X_T\sim p_T^\phi(\cdot|o)}
    [
        r_\beta(o,X_T)
    ]$.
This gives the tilted terminal policy
$\pi_\beta(a|o)=p_T^\beta(a|o)\propto r_\beta(o,a)\pi_{\mathrm{old}}(a|o)$.
Let $u^\beta(x,t,o)$ denote the velocity inducing the marginals $p_t^\beta(\cdot|o)$ and terminal $\pi_\beta(\cdot|o)$.
We define the drift matching objective for flow policies as
\begin{align}
\label{eq:def_DM_fm}
    \mathcal{L}_{\mathrm{DM}}(\theta)
    =
    \E_{t,o, X_t \sim p_t^\beta(\cdot|o)}
    \Big[
        \big\|
            u^\beta(X_t,t,o)
            - u_\theta(X_t,t,o)
        \big\|_2^2
    \Big].
\end{align}
This objective is not implementable since it requires samples from $p_t^\beta$ and access to $u^\beta$.
Therefore, we rewrite \eqref{eq:def_DM_fm} as a conditional drift matching objective that only relies on the behavior process,
\begin{align}
\label{eq:def_cDM_fm}
    \mathcal{L}_{\mathrm{CDM}}(\theta)
    = & ~
    \E_{t,o, X_t \sim p_t^\phi, X_T \sim p_{T| t}^\phi}
    \Big[
        r_\beta(o,X_T)
        \big\|
            u^\beta(X_t| X_T,t,o)
            - u_\theta(X_t,t,o)
        \big\|_2^2
    \Big] \\
    \approx &~
    \E_{t,o, X_0 \sim p_0^\phi, X_T \sim p_T^\phi, X_t \sim p_{t|T,0}^\phi}
    \Big[
        r_\beta(o,X_T)
        \big\|
            \dot{\alpha}_t X_0 + \dot{\beta}_t X_T
            - u_\theta(X_t,t,o)
        \big\|_2^2
    \Big].
    \nonumber
\end{align}
The approximation is reliable when $u_\phi$ is well-trained under the affine path $X_t=\alpha_t X_0+\beta_t X_T$, independent coupling and the trust-region update keeps $\pi_{\mathrm{old}}$ close to the base behavior.
We summarize the training procedure in \cref{alg:rwdm_ppo_fm}.
Finally, we show $\mathcal{L}_{\mathrm{DM}}$ and $\mathcal{L}_{\mathrm{CDM}}$ are equivalent for learning.

\begin{theorem}[Gradient Equivalence for Flow]
\label{thm:grad_equi_fm}
Let $\mathcal{L}_{\mathrm{DM}}$ and $\mathcal{L}_{\mathrm{CDM}}$ be defined in \eqref{eq:def_DM_fm} and \eqref{eq:def_cDM_fm}, respectively.
Fix $o \in\mathcal S$, define normalizing constant
$Z_{\beta}(o) := \E_{X_T \sim p_T^\phi(\cdot|o)}\big[r_\beta(o,X_T)\big]$ and
\begin{align*}
    \ell_{\mathrm{DM}}(\theta,o)
    &:=
    \E_{t, X_t \sim p_t^\beta(\cdot|o)}
    \Big[
        \big\|
            u^\beta(X_t,t,o)
            - u_\theta(X_t,t,o)
        \big\|_2^2
    \Big].
    \\
    \ell_{\mathrm{CDM}}(\theta,o)
    &:=
    \E_{t, X_t \sim p_t^\phi, X_T \sim p_{T| t}^\phi}
    \Big[
        r_\beta(o,X_T)
        \big\|
            u^\beta(X_t| X_T,t,o)
            - u_\theta(X_t,t,o)
        \big\|_2^2
    \Big].
\end{align*}
Then their gradients satisfy
$\nabla_\theta \ell_{\mathrm{CDM}}(\theta,o)
    =
    Z_{\beta}(o) \cdot \nabla_\theta \ell_{\mathrm{DM}}(\theta,o)$.
\end{theorem}

\begin{proof}
Please see \cref{sec:proof_grad_equi_fm} for a detailed proof.
\end{proof}

\textbf{Ratio Regularization.}
In policy gradient methods, the probability ratio $r_\theta(o,a)=\pi_\theta(a|o)/\pi_{\mathrm{old}}(a|o)$ is induced by two normalized policies and therefore satisfies
$\E_{a\sim \pi_{\mathrm{old}}(\cdot|o)}[r_\theta(o,a)]=1$ for each state $o$.
In our setting, $r_\beta(o,a)$ is learned as an auxiliary model, so this normalization constraint is not automatically satisfied.
We therefore add a ratio regularizer that penalizes violations of $\E_{a\sim \pi_{\mathrm{old}}(\cdot|o)}[r_\beta(o,a)]=1$ in \cref{alg:rwdm_ppo,alg:rwdm_ppo_fm}, which stabilizes ratio learning.

\subsection{Unified Framework and Instantiations}
\label{sec:unified_framework}

DoobPO is a unified framework with two design choices.
The first is the generative policy class.
The second is the policy update method that specifies the target policy improvement.
Different choices along these two dimensions induce a broad family of diffusion/flow matching RL algorithms.

Under this view, several prior methods arise as special cases of DoobPO.
In particular, they correspond to policy updates that admit closed-form solutions.
In these cases, the ratio-learning stage becomes unnecessary, and DoobPO reduces to a single reweighted drift matching step.
For example, DPMD \cite{ma2025efficient} corresponds to a policy update with an analytic target ratio.
Under a KL-proximal policy mirror descent update \cite{zhan2023policy}, the updated policy satisfies
$\pi(a|o)\propto \pi_{\mathrm{old}}(a|o)\exp(\hat{A}_{\mathrm{old}}(o,a)/\lambda)$.
Therefore, the corresponding probability ratio satisfies
$r(o,a)\propto \exp(\hat{A}_{\mathrm{old}}(o,a)/\lambda)$.

For more general trust-region policy-gradient methods such as PPO \cite{schulman2017proximal} and GRPO \cite{shao2024deepseekmath}, the target ratio does not admit a closed form.
DoobPO then learns a lightweight ratio model and distills the resulting policy improvement into the generative policy through drift matching.
In this work, we derive DoobPPO and DoobGRPO for both diffusion policies and flow policies by instantiating the RL update with PPO and GRPO.
Hence, the same framework covers both prior closed-form methods and our new learned-ratio instantiations.
We summarize representative examples in \cref{tab:unified_framework}.

\subsection{Optimality from the Perspective of Stochastic Optimal Control}
\label{sec:soc}


\cref{sec:diff_dr_matching} shows that Doob's Drift Matching recovers the ratio-tilted action distribution through Doob's $h$-transform.
In this section, we establish a stronger optimality property from the perspective of endpoint-constrained stochastic optimal control (SOC).
Specifically, among all controlled reverse diffusions
with prescribed endpoint marginals $p_T^h$ and $p_0^h$,
the minimizer of conditional drift matching
induces the minimum-energy drift modification
of the pretrained reverse diffusion process.
Here, minimality is measured by the quadratic control energy.
We now formulate this result.

\begin{proposition}[Conditional Drift Matching Solves the SOC Problem]
\label{prop:cdm_solves_soc}
Let $h_t$ be defined in \eqref{eq:solution_exp}.
Fix a state $o \in\mathcal S$ and let $\mathbb P^\phi(\cdot|o)$ be the reference path measure induced by the reverse-time SDE
\begin{align*}
    \dd X_t
    =
    \bigl[
        f(X_t,t)-\sigma_t^2 s_\phi(X_t,t,o)
    \bigr]\dd t
    +
    \sigma_t\dd W_t,
\end{align*}
where $s_\phi(X_t,t,o)$ is the base model.
Consider the conditional drift matching objective
\begin{align*}
\mathcal{L}_{\mathrm{CDM}}(\theta)
=
\mathbb E_{t, X_t\sim p_t^\phi, X_0\sim p_{0|t}^\phi}
\left[
h_0(X_0,o)
\left\|
\nabla_x\log p_{0|t}^\phi(X_0|X_t,o)
-
s_\theta(X_t,t,o)
\right\|_2^2
\right].
\end{align*}
If the model class $\{s_\theta : \theta \in \Theta\}$ is realizable and the minimum of $L_{\mathrm{CDM}}$ is attained, then any minimizer $s_\theta^\star$ induces the optimal control
$u_t^\star(x,o) := -s_\theta^\star(x,t,o)$
that solves the SOC problem
\begin{align}
    & ~ \quad \quad \quad
    \quad \quad \quad \quad \quad \inf_u
    \frac12
    \mathbb E_{\mathbb Q^u(\cdot|o)}
    \int_0^T
    \|\sigma_t
    u_t(Y_t,o)\|_2^2\dd t,
    \label{eq:oj_soc}
    \\
    \text{s.t.} & ~
    \dd Y_t
    =
    \bigl[
        f(Y_t,t)-\sigma_t^2 s_\phi(Y_t,t,o) + \sigma_t^2 u_t(Y_t,o)
    \bigr]\dd t
    +
    \sigma_t\dd W_t ,
    ~ Y_T \sim p_T^h, ~Y_0 \sim p_0^h,
    \label{eq:sde_soc}
\end{align}
where $\mathbb Q^u(\cdot|o)$ is the path measure induced by the controlled reverse-time SDE in \eqref{eq:sde_soc}, initialized at $t=T$ with $Y_T\sim p_T^h$, and $p_0^h,p_T^h$ are the endpoint distributions given by \eqref{eq:p_t^h}.
\end{proposition}

\begin{proof}
Please see \cref{sec:proof_cdm_learn} for a detailed proof.
\end{proof}

\begin{remark}[Minimum-Energy Modification and Practical Implication]
\label{rem:ratio_h_interpretation}
Given a target terminal tilt, the SOC problem selects, among all controlled reverse diffusions that realize the tilted endpoint marginals, the one with minimum control energy.
The RL methods under our framework are special cases of this principle: we choose $h_0(a,o)\propto r_\beta(o,a)$, so that $p_0^h(a|o)\propto r_\beta(o,a)p_0^\phi(a|o)$.
\cref{prop:cdm_solves_soc} shows that Conditional Drift Matching does not only recover the reward-guided terminal policy; it learns the minimum-energy drift modification that realizes this policy.
In practice, under the drift parametrization $g_\theta$ and the practical loss \eqref{eq:cdm_loss_prac}, the population target satisfies $g_\theta^\star=s_\phi-u_t^\star$, where $u_t^\star(x,o)$ is the optimal SOC control in \cref{prop:cdm_solves_soc}.
Thus Conditional Drift Matching learns the closest feasible modification of the pretrained score model in the SOC sense.
This provides a theoretical explanation for the stability and efficiency of DoobPO.
\end{remark}